\documentclass[11pt]{article}
\usepackage[a4paper,margin=0.95in]{geometry}
\usepackage{amsmath,amssymb,mathtools,booktabs,graphicx,float,hyperref,enumitem,longtable,array,pdflscape}
\hypersetup{hidelinks}
\setlength{\emergencystretch}{2em}
\newcommand{\FigDir}{../figures}
\newcommand{\TabDir}{../tables}
\title{Unified Computational Benchmark for First-Passage Problems\\Luca / Generating-Function Family vs. Time-Domain Recursion Family}
\author{valley-k-small automated report}
\date{\today}
\begin{document}
\maketitle
\begin{abstract}
This document is the single active computational-method comparison report in the repository. 
It keeps the scientific reports separate, but collapses the comparison line into one bilingual benchmark with six workloads. 
The benchmark rule is \emph{practical native-task fairness}: each report is timed on the task it actually needed to complete, rather than on an artificially synchronized full-tail horizon. 
Under the diagnostic protocol the median ratio $t_{\mathrm{time}}/t_{\mathrm{Luca}}$ is 0.0866; 
time recursion wins on ENC-FIXED, ENC-ANY, TT-C1, REF-S0, while Luca/GF wins on RING-1T-paper, TT-LF1. 
Under the curve protocol the median ratio is 0.0671; 
time recursion wins on ENC-FIXED, ENC-ANY, TT-C1, REF-S0, while Luca/GF wins on RING-1T-paper, TT-LF1.
\end{abstract}
\section{Goal and Comparison Rule}
We compare only two families.
\begin{itemize}[leftmargin=1.4em]
\item \textbf{Luca / generating-function family}: closed-form propagators, defect-reduced resolvent recovery, renewal in transform space, and AW/Cauchy-FFT inversion.
\item \textbf{Time-domain recursion family}: exact propagation on the transient state space, including sparse absorbing recursion and absorbed pair-distribution updates.
\end{itemize}
The main benchmark is not a forced full-tail comparison. It is a report-production comparison: how much time does each family need to deliver the quantity that the scientific report actually required?
\section{Unified Notation and Runtime Protocol}
Let $X_t$ be a discrete-time Markov chain, $T$ a target set, $\tau_T=\inf\{t\ge 1:X_t\in T\}$ the first-passage time,
\begin{align}
f_T(t) &= \mathbb P(\tau_T=t),\\
S_T(t) &= \mathbb P(\tau_T>t),\\
\widetilde f_T(z) &= \sum_{t\ge 1} f_T(t) z^t,\\
\widetilde P_{ij}(z) &= \sum_{t\ge 0} \mathbb P_i(X_t=j) z^t = \bigl[(I-zP)^{-1}\bigr]_{ij}.
\end{align}
For single-target problems, renewal closes immediately:
\begin{align}
\widetilde f_{i\to m}(z)=\frac{\widetilde P_{im}(z)}{\widetilde P_{mm}(z)}.
\end{align}
For multiple targets, renewal becomes a small linear system on target-to-target propagators. 
All timings use single-thread BLAS, one warm-up run, three measured runs, and the median wall time as the reported benchmark value.
\section{Luca / Generating-Function Family: Mathematical Spine}
The common transform-domain workflow is:
\begin{enumerate}[leftmargin=1.5em]
\item Build a defect-free propagator or a closed-form single-walker PGF.
\item Recover the defected propagator through a finite-dimensional Woodbury or determinant system.
\item Close a one-target or multi-target renewal relation to obtain the FPT generating function.
\item Invert the PGF numerically with an AW/Cauchy-FFT contour.
\end{enumerate}
The inversion step evaluates
\begin{align}
f(t) = [z^t]\widetilde f(z) = \frac{1}{2\pi i}\oint_{|z|=r} \frac{\widetilde f(z)}{z^{t+1}}\,dz
\approx \frac{r^{-t}}{m}\sum_{k=0}^{m-1}\widetilde f\!\left(re^{2\pi i k/m}\right)e^{-2\pi i kt/m}.
\end{align}
This is the shared numerical shell behind the ring, encounter, two-target, and reflecting GF workloads.
\subsection{Encounter-Specific GF Route}
For the two-walker shortcut encounter, the defect-free pair propagator is
\begin{align}
G_0\big((n_0,m_0)\to(x,y);z\big)=\frac{1}{N^2}\sum_{k_1,k_2}\frac{h_{k_1}(x,n_0)h_{k_2}(y,m_0)}{1-z\lambda_{k_1}^{(1)}\lambda_{k_2}^{(2)}}.
\end{align}
The shortcut becomes a line defect on the pair torus and is restored by
\begin{align}
G(z)=G_0(z)+z\,G_0(z)U\left(I-zV^{\top}G_0(z)U\right)^{-1}V^{\top}G_0(z).
\end{align}
A1 and A8 are therefore not the production solver for the encounter report; they are the spectral anchor from which the new timed GF route is constructed.
\section{Time-Domain Recursion Family: Mathematical Spine}
The time-domain family works directly with absorbing transient updates. If $Q$ is the transient transition matrix and $r$ the transient-to-target flux vector, then
\begin{align}
u_{t+1} &= u_t Q,\\
f(t+1) &= u_t r,\\
S(t+1) &= u_{t+1}\mathbf 1.
\end{align}
For two-target problems, $r$ becomes a two-column target flux block. For the pair encounter, the exact recursion is naturally written in matrix form:
\begin{align}
J_{t+1} &= P_1^{\top}J_tP_2,\\
f_{\mathrm{enc}}(t+1) &= \operatorname{tr}(J_{t+1}),
\end{align}
with the diagonal or designated encounter site zeroed after each step to enforce absorption. 
No asymptotic approximation is used here; the family name refers only to the domain in which the update is carried out.
\section{Workload Instantiations and Configuration Figures}
Each workload is now documented inside this unified report, with its own detailed configuration figure and an explicit formula-level specialization of the two method families.
\begin{table}[H]
\centering\small
\resizebox{\linewidth}{!}{%
\begin{tabular}{lllll}
\toprule
workload & model & states & defects & targets \\
\midrule
RING-1T-paper & ring\_single\_target & 100 & 1 & 1 \\
ENC-FIXED & ring\_encounter\_fixed & 900 & 0 & 1 \\
ENC-ANY & ring\_encounter\_any & 10201 & 101 & 101 \\
TT-C1 & grid2d\_two\_target & 1679 & 632 & 2 \\
TT-LF1 & grid2d\_two\_target & 1679 & 2 & 2 \\
REF-S0 & grid2d\_reflecting & 840 & 9 & 1 \\
\bottomrule
\end{tabular}

}
\caption{The six benchmark workloads kept in the unified comparison.}
\end{table}
\subsection{RING-1T-paper}
This workload is the analytic single-target ring benchmark. The resolvent family starts from the closed-form single-walker propagator
\begin{align}
\widetilde P_{n_0\to n}(u) &= \sum_{k=0}^{N-1} \frac{h_k(n,n_0)}{1-u\alpha_k},\\
\widetilde F_{n_0\to m}(u) &= \frac{\widetilde P_{n_0\to m}(u)}{\widetilde P_{m\to m}(u)},
\end{align}
followed by an AW/Cauchy inversion on a radius-$r$ circle. The time solver instead iterates the absorbing chain directly.
Direct speed mapping for the configuration figure: the diagnostic task is faster with Luca/GF (time=0.007163s, luca=0.0006995s, 10.2x separation). The curve task keeps the same winner (time=0.01769s, luca=0.002326s).
\begin{figure}[p]
\centering
\includegraphics[width=0.97\linewidth]{\FigDir/RING-1T-paper_config_detailed.pdf}
\caption{Detailed ring geometry for the paper-like single-target shortcut benchmark.}
\end{figure}
\clearpage
\subsection{ENC-FIXED}
The fixed-site encounter workload is a single-target first-passage problem on the pair chain. The GF family evaluates the defect-free pair propagator and closes a one-target renewal equation:
\begin{align}
G_0\big((n_0,m_0)\to(x,y);z\big) &= \frac{1}{N^2}\sum_{k_1,k_2}\frac{h_{k_1}(x,n_0)h_{k_2}(y,m_0)}{1-z\lambda^{(1)}_{k_1}\lambda^{(2)}_{k_2}},\\
\widetilde f_{\delta}(z) &= \frac{G\big((n_0,m_0)\to(\delta,\delta);z\big)}{G\big((\delta,\delta)\to(\delta,\delta);z\big)}.
\end{align}
The time family evolves the absorbed pair density exactly on the $N^2$ state space.
Direct speed mapping for the configuration figure: the diagnostic task is faster with Time recursion (time=0.006232s, luca=0.3488s, 56.0x separation). The curve task keeps the same winner (time=0.005959s, luca=0.3508s).
\begin{figure}[p]
\centering
\includegraphics[width=0.97\linewidth]{\FigDir/ENC-FIXED_config_detailed.pdf}
\caption{Detailed configuration for the fixed-site encounter benchmark, including the single pair-state target.}
\end{figure}
\clearpage
\subsection{ENC-ANY}
The anywhere-encounter workload promotes the target from one pair-state to the full diagonal set $D=\{(x,x):x\in\mathbb Z_N\}$. The GF family uses the same defect-free pair propagator but now closes a multi-target renewal system:
\begin{align}
G(z) &= G_0(z) + z\,G_0(z)U\left(I-zV^{\top}G_0(z)U\right)^{-1}V^{\top}G_0(z),\\
\widetilde{\mathbf f}_D(z) &= G_{DD}(z)^{-1}G_{Ds}(z),\\
\widetilde f_{\mathrm{enc}}(z) &= \mathbf 1^{\top}\widetilde{\mathbf f}_D(z).
\end{align}
Here the shortcut acts as a line defect on the pair torus, so the Woodbury correction is the bridge between the defect-free pair spectrum and the actual encounter PGF.
Direct speed mapping for the configuration figure: the diagnostic task is faster with Time recursion (time=0.0123s, luca=2.317s, 188.3x separation). The curve task keeps the same winner (time=0.02134s, luca=4.664s).
\begin{figure}[p]
\centering
\includegraphics[width=0.97\linewidth]{\FigDir/ENC-ANY_config_detailed.pdf}
\caption{Detailed configuration for the anywhere-encounter benchmark, including the diagonal target set and the shortcut-induced defect line.}
\end{figure}
\clearpage
\subsection{TT-C1}
C1 is the practical two-target corridor benchmark. The GF family uses a defect-reduced resolvent instead of a full dense transient solve:
\begin{align}
Q &= Q_0 + U\Delta V^{\top},\\
(I-zQ)^{-1} &= (I-zQ_0)^{-1} + z(I-zQ_0)^{-1}U\left(I-zV^{\top}(I-zQ_0)^{-1}U\Delta\right)^{-1}V^{\top}(I-zQ_0)^{-1},\\
\widetilde{\mathbf f}(z) &= G_{TT}(z)^{-1}G_{Ts}(z),
\end{align}
where $T=\{m_1,m_2\}$. The time solver is the sparse exact absorbing recursion on the transient graph, continued to the long horizon actually needed by the scientific claim.
Direct speed mapping for the configuration figure: the diagnostic task is faster with Time recursion (time=0.01313s, luca=41.19s, 3137.2x separation). The curve task keeps the same winner (time=0.06642s, luca=41.03s).
\begin{figure}[p]
\centering
\includegraphics[width=0.97\linewidth]{\FigDir/TT-C1_config_detailed.pdf}
\caption{Detailed C1 geometry, showing both the full local-bias field and the corridor support that determines the defect-reduced dimension.}
\end{figure}
\clearpage
\subsection{TT-LF1}
TT-LF1 keeps the same two-target mathematics but collapses the defect support to an ultra-sparse local perturbation. This is the in-repo positive anchor where the defect-reduced GF route can beat long-horizon sparse recursion.
Direct speed mapping for the configuration figure: the diagnostic task is faster with Luca/GF (time=0.1901s, luca=0.01725s, 11.0x separation). The curve task keeps the same winner (time=0.1904s, luca=0.01719s).
\begin{figure}[p]
\centering
\includegraphics[width=0.97\linewidth]{\FigDir/TT-LF1_config_detailed.pdf}
\caption{Detailed LF1 sparse-defect geometry used as the positive Luca/GF anchor.}
\end{figure}
\clearpage
\subsection{REF-S0}
REF-S0 is the reflecting low-defect control. The GF family remains a single-target resolvent solve, but here the full AW route is still numerically feasible because the defect count is deliberately kept small:
\begin{align}
\widetilde F_{s\to m}(z) &= \frac{G_{sm}(z)}{G_{mm}(z)},\\
f(t) &= [z^t]\widetilde F_{s\to m}(z).
\end{align}
The time family is the standard exact absorbing recursion on the reflecting lattice.
Direct speed mapping for the configuration figure: the diagnostic task is faster with Time recursion (time=0.006363s, luca=0.04098s, 6.4x separation). The curve task keeps the same winner (time=0.01758s, luca=0.15s).
\begin{figure}[p]
\centering
\includegraphics[width=0.97\linewidth]{\FigDir/REF-S0_config_detailed.pdf}
\caption{Detailed reflecting low-defect S0 control geometry, highlighting the sticky site and the partially permeable barrier.}
\end{figure}
\clearpage
\section{Main Benchmark Results}
\begin{table}[H]
\centering\small
\resizebox{\linewidth}{!}{%
\begin{tabular}{llll}
\toprule
workload family & Luca/GF family & time recursion family & comment \\
\midrule
RING-1T-paper & $O(mN)+O(m\log m)$ & $O(TN)$ & closed-form single-target PGF plus AW \\
ENC-FIXED & $O(mN^3)$ & $O(TN^3)$ & single-target renewal on the pair chain \\
ENC-ANY & $O(mN^3)$ & $O(TN^3)$ & diagonal target set with $|D|=N$ \\
TT-C1 / TT-LF1 & $O(mM^3)+O(m\log m)$ & $O(T|E|)$ & defect-reduced AW vs sparse exact recursion \\
REF-S0 & $O(mn_T^3)+O(m\log m)$ & $O(T|E|)$ & low-defect full AW vs exact recursion \\
\bottomrule
\end{tabular}

}
\caption{Family-level complexity summary.}
\end{table}
\begin{table}[H]
\centering\small
\resizebox{\linewidth}{!}{%
\begin{tabular}{llll}
\toprule
workload & source report & Luca/GF solver & time solver \\
\midrule
RING-1T-paper & ring\_lazy\_flux & ring\_analytic\_aw & ring\_time\_absorption \\
ENC-FIXED & ring\_two\_walker\_encounter\_shortcut & encounter\_fixedsite\_gf\_aw & pair\_fixedsite\_time\_recursion \\
ENC-ANY & ring\_two\_walker\_encounter\_shortcut & encounter\_anywhere\_gf\_aw & pair\_time\_recursion \\
TT-C1 & grid2d\_two\_target\_double\_peak & two\_target\_defect\_reduced\_aw & two\_target\_sparse\_exact \\
TT-LF1 & grid2d\_two\_target\_double\_peak & two\_target\_defect\_reduced\_aw & two\_target\_sparse\_exact \\
REF-S0 & grid2d\_reflecting\_bimodality & reflecting\_full\_aw & reflecting\_exact\_recursion \\
\bottomrule
\end{tabular}

}
\caption{Family-internal solver mapping for the six workloads.}
\end{table}
\begin{figure}[H]
\centering
\includegraphics[width=0.98\linewidth]{\FigDir/unified_runtime_config_overview.pdf}
\caption{One-glance benchmark overview. The numbered mini-panels 1--6 are the six workload configurations, and the numbered groups on the top runtime chart follow the same order. This is the main practical benchmark figure used in the report.}
\end{figure}
\begin{table}[H]
\centering\small
\resizebox{\linewidth}{!}{%
\begin{tabular}{lrrrrl}
\toprule
 workload & Luca/GF (s) & time recursion (s) & effective $T$ (L/T) & recommended & $L^1$ \\
\midrule
RING-1T-paper & 0.0006995 & 0.007163 & 800/800 & Luca / generating-function family & 1.643e-06 \\
ENC-FIXED & 0.3488 & 0.006232 & 900/900 & Time-domain recursion family & 1.382e-10 \\
ENC-ANY & 2.317 & 0.0123 & 420/420 & Time-domain recursion family & 8.994e-14 \\
TT-C1 & 41.19 & 0.01313 & 800/1200 & Time-domain recursion family & 9.912e-10 \\
TT-LF1 & 0.01725 & 0.1901 & 80/12000 & Luca / generating-function family & 2.135e-10 \\
REF-S0 & 0.04098 & 0.006363 & 300/300 & Time-domain recursion family & 1.733e-09 \\
\bottomrule
\end{tabular}

}
\caption{Diagnostic-task runtime table.}
\end{table}
\begin{figure}[H]
\centering
\includegraphics[width=0.98\linewidth]{\FigDir/unified_speedup_by_workload.pdf}
\caption{Per-workload ratio $t_{\mathrm{time}}/t_{\mathrm{Luca}}$. Ratios above one favor the Luca/GF family.}
\end{figure}
\begin{figure}[H]
\centering
\includegraphics[width=0.98\linewidth]{\FigDir/unified_runtime_curve.pdf}
\caption{Curve-task runtime, kept secondary so that long-tail costs do not overwrite the practical report-production benchmark.}
\end{figure}
\begin{table}[H]
\centering\small
\resizebox{\linewidth}{!}{%
\begin{tabular}{lrrrrl}
\toprule
 workload & Luca/GF (s) & time recursion (s) & effective $T$ (L/T) & recommended & $L^1$ \\
\midrule
RING-1T-paper & 0.002149 & 0.01741 & 2000/2000 & Luca / generating-function family & 7.588e-10 \\
ENC-FIXED & 0.3167 & 0.006537 & 900/900 & Time-domain recursion family & 1.382e-10 \\
ENC-ANY & 6.586 & 0.07917 & 800/800 & Time-domain recursion family & 1.045e-13 \\
TT-C1 & 74.17 & 0.06406 & 800/6000 & Time-domain recursion family & 9.912e-10 \\
TT-LF1 & 0.01771 & 0.1882 & 80/12000 & Luca / generating-function family & 2.135e-10 \\
REF-S0 & 0.2829 & 0.01666 & 1200/1200 & Time-domain recursion family & 3.005e-09 \\
\bottomrule
\end{tabular}

}
\caption{Curve-task runtime table.}
\end{table}
\section{Recommendation Matrix}
\begin{table}[H]
\centering\small
\resizebox{\linewidth}{!}{%
\begin{tabular}{lp{0.66\linewidth}}
\toprule
Scenario & Recommendation \\
\midrule
Single-target ring with analytic PGF & Prefer Luca/GF when the closed form already exists and longer windows are needed. \\
Shortcut encounter (anywhere/fixed-site) & Prefer time recursion for production figures; use GF when transform-domain renewal validation is itself the task. \\
2D two-target, medium/high defect & Prefer time recursion; C1-like cases still punish the defect-reduced AW route. \\
2D two-target, ultra-sparse defect & Luca/GF can win; TT-LF1 is the in-repo positive anchor. \\
2D reflecting low-defect control & Time recursion still wins on wall time in this repo. \\
\bottomrule
\end{tabular}

}
\caption{Decision matrix derived from the six-workload benchmark.}
\end{table}
\section{Conclusion}
The comparison line is now fully centralized in one report. The scientific reports remain focused on science, while this document owns the algorithmic comparison, the fairness rule, the workload inventory, and the implementation-level guidance for choosing between transform-domain and time-domain solvers.
\clearpage
\appendix
\section{Appendix A: Symbol Table and Variable Mapping}
\begin{align}
P &: \text{full transition operator}, & Q &: \text{transient transition operator},\\
r &: \text{transient-to-target flux}, & U\Delta V^{\top} &: \text{defect update},\\
G_0(z) &: (I-zQ_0)^{-1}, & G(z) &: (I-zQ)^{-1},\\
m &: \text{AW FFT length}, & r &: \text{AW contour radius}.
\end{align}
\section{Appendix B: Detailed GF Derivations}
The GF family reduces every workload to the same algebraic template. Starting from a defect-free $Q_0$, write
\begin{align}
Q &= Q_0 + U\Delta V^{\top},\\
I-zQ &= (I-zQ_0)\Bigl[I-z(I-zQ_0)^{-1}U\Delta V^{\top}\Bigr],\\
(I-zQ)^{-1} &= (I-zQ_0)^{-1} + z(I-zQ_0)^{-1}U\Bigl[I-zV^{\top}(I-zQ_0)^{-1}U\Delta\Bigr]^{-1}V^{\top}(I-zQ_0)^{-1}.
\end{align}
For a target set $T=\{m_1,\dots,m_K\}$, renewal is written as
\begin{align}
G_{sT}(z)=\widetilde{\mathbf f}_T(z)^{\top}G_{TT}(z),\\
\widetilde{\mathbf f}_T(z)=G_{TT}(z)^{-1}G_{Ts}(z).
\end{align}
The one-target formula is the $K=1$ specialization of this matrix equation.
\section{Appendix C: Detailed Time-Recursion Derivations}
Given the transient state vector $u_t$, exact recursion uses
\begin{align}
u_0 &= e_s,\\
u_{t+1} &= u_tQ,\\
f(t+1) &= u_tr,\\
S(t+1) &= u_{t+1}\mathbf 1. 
\end{align}
For two-target problems, let $R=[r^{(1)}\;r^{(2)}]$ and then $f_1(t+1)=u_tr^{(1)}$, $f_2(t+1)=u_tr^{(2)}$. 
For pair encounters, the matrix recursion form is numerically convenient because the encounter condition is a structured absorbing set on the pair torus.
\section{Appendix D: Per-Workload Formula Specialization}
RING-1T-paper uses a scalar renewal. ENC-FIXED uses a single-target renewal on the pair chain. ENC-ANY uses a diagonal target-set renewal on the pair torus. 
TT-C1 and TT-LF1 use a two-target renewal after defect-reduced selected-propagator recovery. REF-S0 uses the single-target reflecting resolvent. 
These are exactly the formulas instantiated in Section 5 alongside the detailed configuration figures.
\section{Appendix E: Complexity, Memory, Error Metrics, and Benchmark Protocol}
The runtime comparison is accepted only after a same-window numerical agreement check between the two families. 
For each workload, the overlap window is reported through $L^1$, $L^{\infty}$, and a peak-location consistency flag. 
This means wall time is never reported without a paired numerical sanity check.
\section{Appendix F: Embedded Historical Full-FPT Note}
The repository no longer keeps an external active compare report for the old fixed full-tail fairness question. 
Instead, the historical note is embedded here so that the active comparison line remains single-source.
\begin{table}[H]
\centering\small
\resizebox{\linewidth}{!}{%
\begin{tabular}{ll}
\toprule
quantity & value \\
\midrule
diagnostic median $t_{time}/t_{Luca}$ & 0.04899 \\
curve median $t_{time}/t_{Luca}$ & 0.03976 \\
historical policy retained here & embedded appendix only \\
external compare report & retired \\
\bottomrule
\end{tabular}

}
\caption{Embedded historical fairness note retained inside the unified report.}
\end{table}
\end{document}
